\originalchapter{电子结构与原子尺度}
\originalsection{能量尺度与多电子近似}
氢样原子的库仑场给出$a_0=4\pi\epsilon_0\hbar^2/(me^2)$与$E_h=\hbar^2/(ma_0^2)=\alpha^2mc^2$. 因而电子速度量级为$\alpha c$, 精细结构约比束缚能小$\alpha^2$, 超精细结构再含电子核质量比. 这些尺度在实验设计中比背诵单个常数更有用. 本卷计算用SI, 遇到原题使用Gaussian单位时明确注明, 不把$e^2/r$与$e^2/(4\pi\epsilon_0r)$混用.

碱金属的一个价电子远离闭壳层. 在大$r$处它感受到有效电荷$+e$, 但低角动量轨道穿透核芯, 增加束缚. 常用$E_{nl}=-E_h/[2(n-\delta_l)^2]$表述量子亏损$\delta_l$; 该式是对高激发态的有效参数化, 不是把氢公式对所有多电子原子照搬. 高$l$轨道的穿透小, 量子亏损也小.

\originalsection{氦的变分、交换与自旋}
氦的空间哈密顿量为
\[H=\sum_{i=1}^2\left[-\frac{\hbar^2}{2m}\nabla_i^2-\frac{Ze^2}{4\pi\epsilon_0r_i}\right]+\frac{e^2}{4\pi\epsilon_0r_{12}}.\]
以归一化$1s$轨道$\phi_{Z'}=(Z'/a_0)^{3/2}\e^{-Z'r/a_0}/\sqrt\pi$构成乘积, 配反对称自旋单态. 由球对称密度计算库仑积分, 先把$1/|\mathbf r_1-\mathbf r_2|$的角平均化为$1/\max(r_1,r_2)$, 再把径向积分分成$r_1<r_2$和$r_1>r_2$, 得$\langle a_0/r_{12}\rangle=5Z'/8$. 单电子动能为$Z'^2E_h/2$, 核吸引能为$-ZZ'E_h$, 从而
\[E(Z')=E_h\left[Z'^2-2ZZ'+\frac58Z'\right].\]
极小值在$Z'=Z-5/16$, 氦取$Z=2$, 得$E=-729E_h/256\simeq-77.5\,\mathrm{eV}$. 与真实能量的差不是数值积分误差, 而是乘积态未表达$r_{12}$关联.

若两个电子占不同轨道$u,v$, 对称和反对称空间态为$[u(1)v(2)\pm v(1)u(2)]/\sqrt2$. 库仑能分别为$J\pm K$, 其中$J$为直接积分,$K$为交换积分. 交换不是新添一个经典自旋力, 而是费米反对称性改变空间概率. 自旋三重态配空间反对称态, 对同类轨道通常能量更低. 纯电偶极算符不作用于自旋, 在自旋轨道耦合忽略时不能连接单态和三重态; 微弱混合才使禁戒线出现.

\originalsection{精细结构、Lamb位移与有限核大小}
非相对论哈密顿量的最低相对论修正包括$-p^4/(8m^3c^2)$、自旋轨道项和Darwin接触项. 中心势的自旋轨道项为$[2m^2c^2r]^{-1}(\dd V/\dd r)\mathbf L\cdot\mathbf S$, 因$2\mathbf L\cdot\mathbf S=\mathbf J^2-\mathbf L^2-\mathbf S^2$, 本征值由$j,l,s$确定. 氢的三项合并给
\[\Delta E_{nj}=-\frac{mc^2(Z\alpha)^4}{2n^3}\left(\frac1{j+1/2}-\frac3{4n}\right).\]
$2S_{1/2}$与$2P_{1/2}$在这个近似仍简并; 实验的Lamb位移表明辐射自能、真空极化等量子电动力学效应不可省略. 本卷不重推QED重整化, 只区分上述结构修正与辐射修正, 计算涉及Lamb位移时使用原题给定数值.

把点核换成小的球对称电荷分布, 低能展开产生接触修正$\delta V=[Ze^2\langle r_N^2\rangle/(6\epsilon_0)]\delta^3(\mathbf r)$. 因而$\Delta E_{nS}=Ze^2|\psi_{nS}(0)|^2\langle r_N^2\rangle/(6\epsilon_0)$, 正号表示有限核减弱短程吸引. $l>0$波函数在原点为零, 领先接触项消失. 有限大小、精细结构和Lamb位移由不同机制产生, 在精密谱学中必须分别计入.
